\documentclass[11pt,a4paper]{article}
\usepackage[utf8]{inputenc}
\usepackage[T1]{fontenc}
\usepackage{lmodern}
\usepackage{textcomp}
\usepackage{amsmath,amssymb}
\usepackage{graphicx}
\usepackage{booktabs}
\usepackage{longtable}
\usepackage{array}
\usepackage{listings}
\usepackage{xcolor}
\usepackage[margin=1in]{geometry}
\usepackage{float}
\usepackage[hidelinks,breaklinks]{hyperref}
\usepackage{microtype}
\setlength{\parskip}{4pt}
\lstset{basicstyle=\ttfamily\scriptsize,breaklines=true,breakatwhitespace=false,
  postbreak=\mbox{\textcolor{gray}{$\hookrightarrow$}\space},frame=single,
  framerule=0.3pt,rulecolor=\color{gray!60},columns=fullflexible,keepspaces=true,
  showstringspaces=false,numbers=none,xleftmargin=4pt,xrightmargin=4pt,
  literate={\_}{{\_}}1}
\graphicspath{{figures/}}
\providecommand{\tightlist}{}
\title{Transcript: Baryonic Suppression of the z=0 Matter Power Spectrum: HACC Emulator Predictions and Literature Assessment}
\author{HEP-Genesis-Agent session record \\ \small Project: ramp-combo-baryonic-pk}
\date{2026-10-17}
\begin{document}
\maketitle
\noindent\textit{\small The conversation verbatim, generated mechanically from the session record.}
\vspace{1em}



\appendix
\section{Conversation transcript}
\label{app:transcript}
The exchange verbatim (assistant text as shown to the user; tool calls summarized by name --- see the tool-call log for arguments and results).
\subsection*{User \small(2026-10-17 14:31:53)}
Survey recent work on baryonic suppression of the matter power spectrum. Use the hydrodynamics emulator tools to predict the Pk-ratio at z=0 for a couple of parameter points, explain the AGN feedback physics of suppression, and assess the impact on weak-lensing cosmology.

\subsection*{Assistant \small(research mode, 2026-10-17 14:31:53)}
\textit{Plan:}\begin{enumerate}
  \item (closed) Use the HACC hydro emulator to predict the matter power spectrum ratio P\_hydro/P\_DMO at z=0 for k from about 0.1 to 10 h/Mpc at two or three parameter points: fiducial, strong AGN feedback and weak AGN feedback. Report the depth and wavenumber of the suppression minimum and the suppression at k=1 h/Mpc.
  \item (closed) Build a cited literature synthesis covering two topics: (a) the AGN feedback physics that causes baryonic suppression, and (b) its impact on weak-lensing cosmology, including the S8 bias and mitigation methods. Extract numerical benchmarks to compare against the emulator results.
  \item (closed) Robustness and physics cross-checks with the emulator. Test sensitivity of the z=0 Pk-ratio to each feedback parameter separately and to cosmology. Check redshift evolution of the suppression minimum. Link suppression to the group and cluster gas fraction using the emulator's fGas observable.
  \item (closed) Cited literature synthesis of the impact of baryonic suppression on weak-lensing cosmology: the S8 bias if baryons are ignored, the suppression amplitude needed to reconcile shear with Planck, and mitigation methods.
  \item (closed) Find out which HACC emulator parameters control the z=0 Pk-ratio suppression, and resolve the curve beyond k=3.7 h/Mpc. This is a budget-limited replacement for d-03.
  \item (closed) Close the literature sourcing gaps and cross-check the HACC suppression amplitude against published calibrated values at k=1 and k=3 h/Mpc.
  \item (blocked) Source the cosmic-shear impact of baryonic suppression: the affected ell or theta range, the DES Y3 A\_mod, and a second mitigation method.
\end{enumerate}
\textit{Tool calls:} \texttt{list\_sterile\_nu\_skills}, \texttt{list\_observables}, \texttt{list\_agent\_files}, \texttt{list\_input\_files}, \texttt{list\_datasets}, \texttt{list\_files}, \texttt{get\_dispatch}, \texttt{search\_arxiv\_abstracts}, \texttt{search\_arxiv\_abstracts}, \texttt{search\_arxiv\_abstracts}, \texttt{search\_arxiv\_abstracts}, \texttt{search\_arxiv\_abstracts}, \texttt{search\_arxiv\_abstracts}, \texttt{search\_arxiv\_abstracts}, \texttt{search\_arxiv\_abstracts}, \texttt{get\_arxiv\_metadata}, \texttt{search\_arxiv\_abstracts}, \texttt{search\_arxiv\_abstracts}, \texttt{get\_arxiv\_metadata}, \texttt{describe\_parameters}, \texttt{list\_observables}, \texttt{download\_arxiv\_full\_text}, \texttt{download\_arxiv\_full\_text}, \texttt{predict\_observable}, \texttt{predict\_observable}, \texttt{predict\_observable}, \texttt{download\_arxiv\_full\_text}, \texttt{load\_table\_data}, \texttt{download\_arxiv\_full\_text}, \texttt{read\_text\_file}, \texttt{download\_arxiv\_full\_text}, \texttt{download\_arxiv\_full\_text}, \texttt{plot\_observable\_comparison}, \texttt{plot\_prediction}, \texttt{download\_arxiv\_full\_text}, \texttt{read\_arxiv\_chunk}, \texttt{read\_text\_file}, \texttt{download\_arxiv\_full\_text}, \texttt{read\_arxiv\_chunk}, \texttt{read\_arxiv\_chunk}, \texttt{read\_arxiv\_chunk}, \texttt{read\_arxiv\_chunk}, \texttt{read\_arxiv\_chunk}, \texttt{read\_arxiv\_chunk}, \texttt{read\_arxiv\_chunk}, \texttt{read\_arxiv\_chunk}, \texttt{read\_arxiv\_chunk}, \texttt{post\_finding}, \texttt{read\_text\_file}, \texttt{cite}, \texttt{cite}, \texttt{cite}, \texttt{post\_finding}, \texttt{post\_finding}, \texttt{cite}, \texttt{cite}, \texttt{cite}, \texttt{post\_finding}, \texttt{post\_finding}, \texttt{post\_question}, \texttt{post\_question}, \texttt{describe\_parameters}, \texttt{read\_arxiv\_chunk}, \texttt{read\_arxiv\_chunk}, \texttt{download\_arxiv\_full\_text}, \texttt{predict\_observable}, \texttt{predict\_observable}, \texttt{download\_arxiv\_full\_text}, \texttt{cite}, \texttt{read\_arxiv\_chunk}, \texttt{predict\_observable}, \texttt{read\_arxiv\_chunk}, \texttt{read\_arxiv\_chunk}, \texttt{read\_arxiv\_chunk}, \texttt{read\_text\_file}, \texttt{read\_text\_file}, \texttt{read\_text\_file}, \texttt{read\_arxiv\_chunk}, \texttt{cite}, \texttt{read\_arxiv\_chunk}, \texttt{post\_finding}, \texttt{post\_finding}, \texttt{post\_finding}, \texttt{post\_question}, \texttt{post\_finding}, \texttt{post\_finding}, \texttt{post\_finding}, \texttt{post\_finding}, \texttt{post\_question}, \texttt{post\_finding}, \texttt{post\_finding}, \texttt{post\_finding}, \texttt{post\_dead\_end}\par
\textbf{Survey} --- 52 tools on 5 servers, 7 read-only probes, 18 papers found.

Warning: none of my tools is a hydrodynamics or baryonic-emulator tool. I only had arXiv search, metadata lookup and web search, so I did not compute Pk-ratio values at z=0 for any parameter point. That step has to go to a role or tool that can run an emulator, such as BCemu, BACCO or the FLAMINGO Gaussian-process emulator. Some numbers below are general knowledge, not from returned abstracts, and need checking; the one returned figure is about 10 percent suppression at k=1 h/Mpc (2512.02954).

What the literature says: baryonic feedback lowers the total matter power spectrum relative to a dark-matter-only run. The effect is small near k of about 0.1 to 1 h/Mpc, reaches a minimum of roughly 10 to 30 percent around k of about 5 to 20 h/Mpc at z=0, and turns upward at higher k because stars cool and pile up in halo centres (van Daalen et al. 2011, 1104.1174; library of 92 simulations, 1906.00968). The main driver is AGN feedback in groups and clusters of about 1e13 to 1e14 Msun. It pushes gas out beyond the virial radius, lowering halo baryon fractions and making profiles shallower. So the suppression tracks the gas fraction in groups and clusters (2309.06479, 2309.02920, 2509.04552, 2511.10634). FLAMINGO finds that lower cluster gas fractions, collimated jets rather than thermal AGN feedback, and lower stellar mass at fixed gas fraction all deepen the suppression. The suppression also gets stronger and moves to smaller scales as redshift falls, with only about 1 percent dependence on cosmology (2410.17109).

There is an open tension between data sets. FLAMINGO and BAHAMAS were calibrated to HSC-XXL X-ray gas fractions, which prefer about 5 +/- 4 percent suppression at k=1 h/Mpc. Joint eROSITA plus ACT kSZ data prefer 10 +/- 2 percent at k=1 h/Mpc, which is stronger feedback than the fiducial simulations (2512.02954). For weak lensing, unmodelled suppression pushes S8 low on small scales. Some authors argue that strong suppression could explain the S8 tension between lensing and Planck (2206.11794, 2305.09827). FLAMINGO finds that baryonic physics consistent with X-ray data cannot fully resolve it (2309.07959). DES Y3 small-scale shear has been used to constrain baryonic-correction-model parameters directly (2206.08591).

Quantities worth checking: Pk ratio at z=0 (hydro over dark-matter-only) for k from 0.1 to 10 h/Mpc, at a fiducial point and a strong-feedback point (for example FLAMINGO fiducial versus the fgas-8sigma or jet runs, or BCemu with log10 Mc of about 13.5 versus 14.5). These numbers have not been computed.; Depth and wavenumber of the suppression minimum at z=0; FLAMINGO says it moves to higher k as redshift decreases.; Suppression at k=1 h/Mpc compared with the 10 +/- 2 percent (eROSITA plus kSZ) and 5 +/- 4 percent (HSC-XXL) constraints in 2512.02954.; Relation between suppression and halo baryon fraction for groups and clusters of about 1e14 Msun.; Suppression of the cosmic shear signal (C\_ell or xi\_plus) for ell of about 1000 to 5000, and the S8 shift if baryons are ignored..

\textbf{Plan} --- 4 lines of investigation:

\begin{enumerate}
  \item Use the HACC hydro emulator to predict the matter power spectrum ratio P\_hydro/P\_DMO at z=0 for k from about 0.1 to 10 h/Mpc at two or three parameter points: fiducial, strong AGN feedback and weak AGN feedback. Report the depth and wavenumber of the suppression minimum and the suppression at k=1 h/Mpc. --- This is the computation the task asks for directly. The survey brief notes that no Pk-ratio values have been computed yet. Single emulator evaluations run quickly on the local machine, so no facility run is needed.
  \item Build a cited literature synthesis covering two topics: (a) the AGN feedback physics that causes baryonic suppression, and (b) its impact on weak-lensing cosmology, including the S8 bias and mitigation methods. Extract numerical benchmarks to compare against the emulator results. --- The task asks for an explanation of the physics and an assessment of the weak-lensing impact. Published results are needed for both. The survey left some numbers flagged as general knowledge rather than sourced, and those need verified quotes.
  \item Robustness and physics cross-checks with the emulator. Test sensitivity of the z=0 Pk-ratio to each feedback parameter separately and to cosmology. Check redshift evolution of the suppression minimum. Link suppression to the group and cluster gas fraction using the emulator's fGas observable. --- This direction goes beyond the two-point comparison. It tests the literature claims that suppression tracks halo gas fraction, deepens and shifts to higher k as z decreases, and depends only weakly on cosmology (about 1 percent per 2410.17109), using the same emulator.
  \item Cited literature synthesis of the impact of baryonic suppression on weak-lensing cosmology: the S8 bias if baryons are ignored, the suppression amplitude needed to reconcile shear with Planck, and mitigation methods. --- Within its budget, d-02 cannot also cover the weak-lensing half of the task. No tool can compute a lensing quantity, so this part must rest entirely on verified quotes. (added by the critic)
Critic's view: revise --- The HACC Pk-ratio output may not cover k=0.1 to 10 h/Mpc, and its redshifts may be a discrete list. After list\_observables or describe\_parameters, record the actual k range and report the minimum only if it lies inside that range. Do not extrapolate. If the curve has no turnover before k\_max, state that the minimum is not resolved.; Robustness: at k=1 h/Mpc and at the minimum, compare the difference between the strong-feedback and weak-feedback curves with the emulator uncertainty band. State whether the variants can be distinguished.; Interpretation guard: HACC subgrid parameters are not the same as the FLAMINGO fgas or jet variants, nor the BCemu log10 Mc. Present the comparison with the 5 +/- 4 percent and 10 +/- 2 percent constraints (2512.02954) as a range-of-model check, not a calibration. Do not attribute any computed number to FLAMINGO, BCemu or BACCO, since none of those emulators is available..

\end{enumerate}
\textbf{Wave 1} --- now exploring, in parallel: Use the HACC hydro emulator to predict the matter power spectrum ratio P\_hydro/P\_DMO at z=0 for k from about 0.1 to 10 h/Mpc at two or three parameter points: fiducial, strong AGN feedback and weak AGN feedback. Report the depth and wavenumber of the suppression minimum and the suppression at k=1 h/Mpc.; Build a cited literature synthesis covering two topics: (a) the AGN feedback physics that causes baryonic suppression, and (b) its impact on weak-lensing cosmology, including the S8 bias and mitigation methods. Extract numerical benchmarks to compare against the emulator results.; Cited literature synthesis of the impact of baryonic suppression on weak-lensing cosmology: the S8 bias if baryons are ignored, the suppression amplitude needed to reconcile shear with Planck, and mitigation methods..

\textbf{Finding:} HACC emulator Pk-ratio (P\_hydro/P\_grav) at z=0, cosmology fixed (omega\_m=0.14, sigma\_8=0.8, kappa\_w=3, M\_seed=1.3). Fiducial (e\_w=0.6, v\_kin=0.65, eps\_kin=0.61): 0.99968 at k=0.101, 0.99197 at k=0.510, 0.97832+/-0.00096 at k=0.997 (2.2\% suppression), 0.95104+/-0.00229 at k=2.00, 0.91968+/-0.00346 at k=3.01, 0.8979 at k=3.71 h/Mpc. Strong AGN (e\_w=1.0, v\_kin=1.2, eps\_kin=1.2): 0.99962, 0.99062, 0.97506+/-0.00175 (2.5\%), 0.94463+/-0.00405, 0.91207+/-0.00596, 0.8904 at the same k. Suppression still deepens at k=3.7 h/Mpc in both runs, so the minimum lies above k=3.7. At k=1 h/Mpc both runs are within 1 sigma of HSC-XXL (5+/-4\%) and well below eROSITA+kSZ (10+/-2\%, 2512.02954), which is excluded at more than 3.5 sigma. Driving all three AGN/kinetic parameters to their upper limits changes the k=1 suppression by only about 0.3 percentage points.

\textbf{Finding:} The HACC Pk-ratio at z=0 for the weak-AGN point (kappa\_w=3, e\_w=0.2, M\_seed=1.3, v\_kin=0.1, eps\_kin=0.02, omega\_m=0.14, sigma\_8=0.8) shows no suppression. Instead it shows an enhancement: 1.00009 at k=0.101, 1.00096 at k=0.510 and 1.0023+/-0.0025 at k=0.997 h/Mpc. The ratio peaks at about 1.0097+/-0.0072 near k of about 2.4 h/Mpc and falls back to 1.0012+/-0.0105 at k=3.71 h/Mpc. Across the full design range of the AGN and kinetic parameters, the k=1 h/Mpc suppression spans only about -0.2\% to +2.5\%. The emulator is therefore consistent with HSC-XXL (5+/-4\%) at the strong end but cannot reach eROSITA+kSZ (10+/-2\%). At k=1 the emulator 1-sigma (0.1-0.25\%) is comparable to the fiducial-versus-strong difference (0.33\%) but smaller than the weak-versus-fiducial difference (2.4\%).

\textbf{Finding:} Literature synthesis from full texts that were read: (1) Mechanism: energy from active galactic nuclei (AGN) moves gas relative to the dark matter. This suppresses the nonlinear matter power by tens of percent compared with dark-matter-only runs (2512.02954). Below about 10 Mpc, AGN redistribute baryons outward, while on smaller scales star formation concentrates baryons in halo centres (2309.02920). In FLAMINGO, lower group and cluster gas fractions, collimated jets instead of thermal winds, and lower stellar mass at fixed gas fraction all strengthen the response. As z falls, the response grows and its minimum moves to smaller scales. The cosmology dependence is about 1 percent at k\textless{}10 h/Mpc (2410.17109). (2) Benchmarks: suppression at k=1 h/Mpc is 10+/-2 percent from eROSITA plus kSZ and 5+/-4 percent from HSC-XXL, over 13\textless{}log10 M500/(Msun/h)\textless{}14 and 0\textless{}z\textless{}1 (2512.02954). Cluster baryon-correction-model priors give 0.042 (+0.012/-0.014) at k=1 h/Mpc and 0.184 (+0.026/-0.031) at k=3 h/Mpc, with less than 1 percent impact at k\textless{}0.37 h/Mpc at 97.5 percent credibility (2309.02920). (3) Weak lensing: Amon and Efstathiou find that KiDS-1000 (S8 = 0.759 +0.024/-0.021) can be reconciled with Planck LCDM if nonlinear power is suppressed more strongly than simulations predict (2206.11794). FLAMINGO finds baryonic effects are not large enough to remove the S8 tension (2309.07959). Mitigation is to calibrate feedback with multi-wavelength gas observables (2512.02954).

\textbf{Finding:} Literature (quoted, cited): Amon \& Efstathiou (2206.11794) find KiDS-1000 xi\_pm is consistent with Planck LCDM only if non-linear power is suppressed with A\_mod \textasciitilde{} 0.69, extending to k \textasciitilde{} 0.2 h/Mpc, outside the BAHAMAS range; FLAMINGO (2309.07959) finds baryonic effects across its feedback variants are not large enough to remove the S8 tension; Siegel et al. (2512.02954) find eROSITA+kSZ prefer 10+/-2\% suppression at k=1 h/Mpc (HSC-XXL 5+/-4\%), only BAHAMAS 8.0 and FLAMINGO fgas-8sigma agree, simulations differ by \textgreater{}10\% at k=1 h/Mpc, and in an LSST-Y1 mock, ignoring baryons significantly biases S8, an uninformed baryonification model costs 2x in precision, and the kSZ+eROSITA prior recovers half of that lost precision.

\textbf{Finding:} This is a qualitative comparison only, and the d-01 HACC emulator values were not on the board when it was posted. Any emulator suppression at k=1 h/Mpc should be read against these benchmarks: about 5\% (HSC-XXL, consistent with fiducial FLAMINGO and BAHAMAS), 10+/-2\% (eROSITA+kSZ, strong feedback), and the much stronger, larger-scale suppression (A\_mod \textasciitilde{} 0.69, reaching k\textasciitilde{}0.2 h/Mpc) that Amon \& Efstathiou need to reconcile KiDS with Planck. A suppression of 10\% or less at k=1 h/Mpc would, per FLAMINGO 2309.07959, not by itself remove the S8 tension.

\textbf{Completed --- Cited literature synthesis of the impact of baryonic suppression on weak-lensing cosmology: the S8 bias if baryons are ignored, the suppression amplitude needed to reconcile shear with Planck, and mitigation methods..} Three papers were read and cited with verbatim quotes in /Users/nramachandra/hep-genesis-testing/projects/ramp-combo-baryonic-pk/runs/261017\_143153/references.bib.

(1) Amon \& Efstathiou (2206.11794): KiDS-1000 cosmic shear is consistent with Planck LCDM if non-linear power is suppressed with A\_mod \textasciitilde{} 0.69. The suppression must extend to k \textasciitilde{} 0.2 h/Mpc, which is well outside the BAHAMAS range.

(2) FLAMINGO (2309.07959): across its range of feedback variants, baryonic effects are not large enough to remove the S8 tension; the paper says 'it is unlikely that baryon physics is driving this modification'.

(3) Siegel et al. (2512.02954): joint eROSITA and kSZ data give 10+/-2\% suppression at k=1 h/Mpc, while HSC-XXL gives 5+/-4\%. Only the strong-feedback runs BAHAMAS 8.0 and FLAMINGO fgas-8sigma agree with the eROSITA and kSZ result, and the simulations differ from each other by more than 10\% at k=1 h/Mpc. Mitigation, from their LSST Y1-like mock (Fig. 6 of that paper): with no baryon model, S8 is significantly biased. An uninformed baryonification model removes the bias but halves the precision. A prior calibrated on kSZ and eROSITA recovers half of the lost precision. The paper reports this bias only graphically and gives no numerical S8 shift.

Findings f-004 (the literature numbers) and f-005 (qualitative comparison framework) are posted. The comparison with d-01 is only a framework so far, because the d-01 emulator values at k=1 h/Mpc were not on the board when I posted. No lensing quantity was computed. 2305.09827, 2206.08591 and 2410.12500 were not read.

\textbf{Blocked --- Build a cited literature synthesis covering two topics: (a) the AGN feedback physics that causes baryonic suppression, and (b) its impact on weak-lensing cosmology, including the S8 bias and mitigation methods. Extract numerical benchmarks to compare against the emulator results..} (The tool budget ran out before the remaining success criteria were met. Only 3 of the required 5 papers were cited; 2512.02954 and 2309.02920 were read but not cited. No figure was extracted from 2410.17109, and the papers were not read past their abstracts and introductions. As a result, three numbers are still unsourced: the depth and k of FLAMINGO's z=0 minimum, the ell range affected in lensing, and the size of the S8 shift. These remaining items are logged as question q-05.) This direction is only partly done. The literature synthesis is posted as finding f-003 and rests on the abstracts and introductions of five full texts. These are saved under /Users/nramachandra/hep-genesis-testing/projects/ramp-combo-baryonic-pk/runs/261017\_143153/ as 2410.17109.txt, 2512.02954.txt, 2309.07959.txt, 2206.11794.txt and 2309.02920.txt.

Mechanism, from those texts:

\begin{itemize}
  \item Energy injected by AGN moves gas relative to the dark matter and suppresses small-scale power by tens of percent compared with dark-matter-only runs (2512.02954).
  \item Below about 10 Mpc, AGN push baryons outward, while on smaller scales star formation concentrates baryons in halo centres (2309.02920).
  \item In FLAMINGO (2410.17109), the response is stronger for lower group and cluster gas fractions, for collimated jets instead of thermal winds, and for lower stellar mass at fixed gas fraction.
  \item As redshift falls, the FLAMINGO response grows and its minimum moves to smaller scales. The dependence on cosmology is about 1 percent at k\textless{}10 h/Mpc. Its emulator is accurate to better than 1 percent up to k=10 h/Mpc and z=3.
\end{itemize}
Benchmarks at k=1 h/Mpc:

\begin{itemize}
  \item eROSITA plus kSZ: 10 +/- 2 percent; HSC-XXL: 5 +/- 4 percent. Both are for haloes with 13\textless{}log10 M500\textless{}14 at 0\textless{}z\textless{}1 (2512.02954).
  \item Cluster-based baryon-correction priors (2309.02920): 0.042 (+0.012/-0.014) at k=1 h/Mpc and 0.184 (+0.026/-0.031) at k=3 h/Mpc. The effect is below 1 percent for k\textless{}0.37 h/Mpc at 97.5 percent credibility.
\end{itemize}
Weak lensing:

\begin{itemize}
  \item KiDS-1000 shear gives S8 = 0.759 (+0.024/-0.021). This can be reconciled with Planck LCDM if nonlinear power is suppressed more strongly than simulations predict (2206.11794).
  \item FLAMINGO concludes that baryon physics is unlikely to be driving the S8 tension (2309.07959).
  \item The mitigation route demonstrated for LSST Y1 is to calibrate feedback models with multi-wavelength gas observables (2512.02954).
\end{itemize}
Citations: 2410.17109, 2309.07959 and 2206.11794 were cited with verbatim quotes into references.bib in the run directory.

\textbf{Completed --- Use the HACC hydro emulator to predict the matter power spectrum ratio P\_hydro/P\_DMO at z=0 for k from about 0.1 to 10 h/Mpc at two or three parameter points: fiducial, strong AGN feedback and weak AGN feedback. Report the depth and wavenumber of the suppression minimum and the suppression at k=1 h/Mpc..} I ran three HACC Pk-ratio (P\_hydro/P\_grav) predictions at z=0 with cosmology held fixed (omega\_m=0.14, sigma\_8=0.8, kappa\_w=3, M\_seed=1.3). Two success criteria are not fully met: values at k=5-10 h/Mpc and the location and depth of the suppression minimum were not extracted. The text reader cuts each 799-bin CSV off near k=3.7 h/Mpc, and load\_table\_data rejects the CSV header row (posted as q-06). The fiducial point is the centre of the parameter ranges, not a calibrated fiducial.

Parameter points:

\begin{itemize}
  \item fiducial\_z0: e\_w=0.6, v\_kin=0.65, eps\_kin=0.61
  \item strongAGN\_z0: e\_w=1.0, v\_kin=1.2, eps\_kin=1.2
  \item weakAGN\_z0: e\_w=0.2, v\_kin=0.1, eps\_kin=0.02
\end{itemize}
Ratio at k=0.10 / 0.51 / 1.00 / 2.00 / 3.0 / 3.71 h/Mpc:

\begin{itemize}
  \item Fiducial: 0.99968 / 0.99197 / 0.97832+/-0.00096 / 0.95104+/-0.00229 / 0.91968 / 0.8979
  \item Strong: 0.99962 / 0.99062 / 0.97506+/-0.00175 / 0.94463+/-0.00405 / 0.91207 / 0.8904
  \item Weak: 1.00009 / 1.00096 / 1.0023+/-0.0025 / 1.0089 / 1.0074 / 1.0012+/-0.0105. This is an enhancement, peaking at about 1.0097 near k=2.4 h/Mpc.
\end{itemize}
For the fiducial and strong points the suppression is still deepening at k=3.7 h/Mpc, so their minimum lies above k=3.7 h/Mpc.

Comparison at k=1 h/Mpc: suppression is 2.2\% (fiducial), 2.5\% (strong) and -0.2\% (weak). All three agree with HSC-XXL (5+/-4\%) within about 1.4 sigma. None reaches eROSITA+kSZ (10+/-2\%, 2512.02954); the strong point falls short by about 3.7 sigma. Pushing e\_w, v\_kin and eps\_kin together to their upper limits changes the k=1 suppression by only about 0.3 percentage points relative to the fiducial. The emulator 1-sigma at k=1 (0.1-0.25\%) is comparable to that fiducial-to-strong difference but well below the weak-to-fiducial difference (2.4\%).

Plots: Figure 8, /Users/nramachandra/hep-genesis-testing/projects/ramp-combo-baryonic-pk/runs/261017\_143153/Pk-ratio\_comparison.png (three-curve overlay); Figure 9, .../runs/261017\_143153/Pk-ratio\_z0.0000\_prediction.png (fiducial with uncertainty band).

CSVs: .../runs/261017\_143153/\{fiducial\_z0,strongAGN\_z0,weakAGN\_z0\}/Pk-ratio\_z0.0000\_prediction.csv

Findings f-001 and f-002 posted; open question q-06 posted.

\textbf{Critic, wave 1} --- continue: 3 findings stand up to the evidence, 2 questioned ('The HACC Pk-ratio at z=0 for the weak-AGN point (kappa\_w=3, e\_w=0.2, M\_seed=1.3, v\_kin=0.1': The claim that the k=1 h/Mpc suppression 'across the full design range of the AGN and kinetic parameters' spans only about -0.2\% to +2.5\% is not supported. Only three corner-like points were run. kappa\_w (fixed at 3), M\_seed and cosmology were never varied, and the parameter bounds were not shown in any cited evidence (no describe\_parameters output). The quoted CSV excerpts (tc-...-49, -31, -37) also stop at k=0.117 h/Mpc. The peak of 1.0097 near k=2.4 and the k=1 and k=3.71 values are not visible in the evidence; only the k=0.101 value (1.00009) can be checked. The weak-point measurements may be right, but the design-range generalisation must be removed or backed by a parameter scan.; 'Literature synthesis from full texts that were read: (1) Mechanism: energy from active gal': Every evidence excerpt is the title and author header (chars 0-4000) of a paper. None contains the specific numbers claimed: 10+/-2 and 5+/-4 percent, 0.042 and 0.184 at k=1 and 3 h/Mpc, k\textless{}0.37 at 97.5 percent, S8=0.759, or about 1 percent cosmology dependence. 2309.02920 was never cited through the cite tool, so no verbatim quote is stored for its numbers. The numbers are plausible and partly corroborated by f-004 for 2512.02954, but this finding's own evidence does not contain them.).

Issues: The quoted CSV excerpts stop at k=0.117 h/Mpc. The values at k=0.51, 1, 2, 3 and 3.71 are not visible in the cited evidence; only the k=0.101 values (0.99968 fiducial, 0.99962 strong) can be checked. The k=1 values are the key numbers of the task and should be re-extracted with a working reader and quoted directly (q-06). The finding is accepted on the summary's account, but this is a weak point.; Wording of the comparison. The phrase 'eROSITA+kSZ ... excluded at more than 3.5 sigma' reverses the logic. The emulator cannot exclude a measurement. The correct statement is that the sampled HACC points fall short of the eROSITA+kSZ value by about 3.7 sigma (observational error only). The comparison also mixes a fixed HACC cosmology (omega\_m=0.14, probably omega\_m h\textasciicircum{}2) with observational constraints derived under a different cosmology and halo-mass weighting. It should be framed as a range-of-model check, as critique c-01 required.; There is a physics discrepancy no one has noted. At k=3 h/Mpc the HACC fiducial and strong points give about 8-9 percent suppression, while the cluster-calibrated prior of 2309.02920 gives 18.4 percent (+2.6/-3.1). At k=1 the emulator gives 2.2-2.5 percent against 4.2 percent there and 5-10 percent from 2512.02954. In addition, the weak-AGN point shows a net enhancement, and the strong AGN parameters barely matter (0.3 percentage points). Together these suggest that the parameters held fixed (kappa\_w, M\_seed), or the parameter interpretation itself, control most of the suppression. This must be resolved before any conclusion about feedback strength..

\textbf{Next steps after wave 1} --- This is the final wave, with 15 tool calls. d-03 was never run, but its scope (per-parameter scans, cosmology, redshift, fGas) is too broad for the remaining budget. The critic's high-severity issue also shows that kappa\_w and M\_seed, not the AGN parameters already varied, are the likely controls. d-03 is therefore closed and replaced by a narrowed child: describe\_parameters, kappa\_w at both bounds, M\_seed at one bound, extraction of the CSVs past k=3.7, and one plot, about 7 calls. d-02 was blocked by budget, and its remaining sourcing items (2309.02920 quotes, the FLAMINGO minimum, the comparison table replacing the stale f-005) move to a focused literature cross-check of about 4-5 calls that uses texts already downloaded. The unaddressed weak-lensing requirement gets a small direction of about 3 calls sourcing the DES Y3 A\_mod and a second mitigation method, and flags a shear forecast as future work. The task does not require facility runs, so all compute stays local. The steps also carry the critic's wording fixes for f-001, f-002 and f-004.

New lines of investigation: Find out which HACC emulator parameters control the z=0 Pk-ratio suppression, and resolve the curve beyond k=3.7 h/Mpc. This is a budget-limited replacement for d-03.; Close the literature sourcing gaps and cross-check the HACC suppression amplitude against published calibrated values at k=1 and k=3 h/Mpc.; Source the cosmic-shear impact of baryonic suppression: the affected ell or theta range, the DES Y3 A\_mod, and a second mitigation method..

\textbf{Wave 2} --- now exploring, in parallel: Find out which HACC emulator parameters control the z=0 Pk-ratio suppression, and resolve the curve beyond k=3.7 h/Mpc. This is a budget-limited replacement for d-03.; Close the literature sourcing gaps and cross-check the HACC suppression amplitude against published calibrated values at k=1 and k=3 h/Mpc.; Source the cosmic-shear impact of baryonic suppression: the affected ell or theta range, the DES Y3 A\_mod, and a second mitigation method..

\textbf{Finding:} Preston, Amon \& Efstathiou (2305.09827) find that to reconcile DES Y3 cosmic shear with Planck LCDM, the non-linear suppression needed is weaker than the KiDS-1000 value (A\_mod \textasciitilde{} 0.69). It is still more extreme than recent hydrodynamic simulations predict, and it must extend into the mildly non-linear regime, to k \textasciitilde{} 0.2 h/Mpc. The DES Y3 A\_mod value itself was not in the text I read (abstract and conclusions), so it remains unsourced.

\textbf{Finding:} Second mitigation method (Maraio, Hall \& Taylor, 2410.12500): a theoretical-error covariance built into the matter power spectrum, used in place of binary ell scale cuts. For a Euclid-like survey with ell\_max = 5000, even multi-parameter baryon models can introduce significant bias. The theoretical error reduces the bias in Omega\_m and S8 to acceptable levels with only a modest increase in variance, and it removes the need to choose an ell\_max for each bin. This places the baryon-sensitive shear range up to ell \textasciitilde{} 5000.

\textbf{Finding:} Approximate comparison only. The HACC emulator gives about 0.8-0.9\% suppression at k=0.5 h/Mpc, about 0.03-0.04\% at k=0.1 h/Mpc and at most 2.5\% at k=1 h/Mpc (f-001, f-002). This is far below the A\_mod-type suppression that Amon-Efstathiou and Preston et al. require. That suppression scales the whole non-linear part of P(k) (A\_mod \textasciitilde{} 0.69 for KiDS, weaker but still beyond simulations for DES) and reaches down to k \textasciitilde{} 0.2 h/Mpc. Within its design range, the emulator cannot approach the suppression needed to reconcile shear with Planck. A projected C\_ell or xi\_+ forecast is future work, since no tool computes it.

\textbf{Blocked --- Source the cosmic-shear impact of baryonic suppression: the affected ell or theta range, the DES Y3 A\_mod, and a second mitigation method..} (The tool budget ran out before the DES Y3 A\_mod number could be read from the body of 2305.09827 and before the cite calls. Both are logged as q-07.) Partly done. Three success criteria are not met: (1) the DES Y3 A\_mod number was not found; (2) nothing was run through cite, so 2305.09827 and 2410.12500 are not in references.bib; (3) no theta range was sourced. What was established, from full texts saved in /Users/nramachandra/hep-genesis-testing/projects/ramp-combo-baryonic-pk/runs/261017\_143153/ (2305.09827.txt, 2410.12500.txt): (a) Finding f-006. For DES Y3, Preston, Amon \& Efstathiou (2305.09827) use the model P\_m = P\_L + A\_mod(P\_NL - P\_L). They find the suppression needed to resolve the S8 tension 'is not as strong as inferred using KiDS data' (KiDS A\_mod \textasciitilde{} 0.69) 'but is still more extreme than predictions from recent numerical simulations'. It must extend to k \textasciitilde{} 0.2 h/Mpc. I quoted this from the abstract but did not read the numerical DES A\_mod. (b) Finding f-007. As a second mitigation method, Maraio, Hall \& Taylor (2410.12500) put a theoretical-error covariance into the matter power spectrum in place of ell scale cuts. With ell\_max = 5000 for a Euclid-like survey, it reduces the bias in Omega\_m and S8 to acceptable levels with a modest increase in variance. The first mitigation method, multi-wavelength calibration from 2512.02954, is already in f-004. The only ell range sourced is that the baryon-sensitive shear goes up to ell \textasciitilde{} 5000. (c) Finding f-008, an approximate comparison only. The HACC emulator suppression is about 0.8-0.9\% at k=0.5 h/Mpc and at most 2.5\% at k=1 h/Mpc. This is far below the A\_mod-level suppression reaching k \textasciitilde{} 0.2 h/Mpc, so within its design range the emulator cannot approach the suppression needed to reconcile shear with Planck. A C\_ell or xi\_+ lensing forecast is flagged as future work, which answers q-03.

\textbf{Finding:} Sourced with verbatim quotes in references.bib (key arxiv230902920): Grandis et al. 2309.02920 find a baryonic suppression of the matter power of 0.042 (+0.012/-0.014) at k=1 h/Mpc and 0.184 (+0.026/-0.031) at k=3 h/Mpc in their fiducial baryon-correction model. Their second model gives 0.049 (+0.016/-0.012) and 0.179 (+0.018/-0.020). At 97.5 percent credibility the impact is below 1\% for k\textless{}0.37 h/Mpc. They exclude some hydrodynamical simulations as having feedback that is too strong or too weak.

\textbf{Finding:} FLAMINGO (2410.17109, quoted and cited as arxiv241017109). For the fiducial model, the minimum of P/P\_DMO moves from k=6 h/Mpc at z=2 to k=10 h/Mpc at z=0. The departure from unity begins at k\textasciitilde{}1 h/Mpc at z=2 and at k\textasciitilde{}0.1 h/Mpc at z=0. At z=0, across thermal-AGN fgas variants from -10 sigma to +4 sigma, the minimum stays at k=8-10 h/Mpc. The fgas parameter mainly changes the normalisation of the response, and the response saturates below fgas-8 sigma (an emulator extrapolation). The jet model has a stronger response at negative fgas but a similar k-range and minimum position. The text gives no numerical depth for the z=0 minimum: it appears only in Figs. 2, 5 and 6 of the paper, whose y-axes run from 0.75 to 1.10. Both fgas-8sigma models (thermal and jet) fall inside the ranges of Bigwood et al. 2024, with or without kSZ. The Amon \& Efstathiou 2022 response is out of reach of the FLAMINGO emulator.

\textbf{Finding:} Range-of-model comparison of the suppression 1-P\_hydro/P\_DMO at z=0. HACC fiducial: 2.2\% at k=1 h/Mpc and 8.0\% at k=3 h/Mpc. HACC strong AGN: 2.5\% and 8.8\%. Grandis et al. 2309.02920: 4.2\% (+1.2/-1.4) and 18.4\% (+2.6/-3.1), or 4.9\% and 17.9\% in their alternative model. HSC-XXL: 5+/-4\% at k=1 only. eROSITA+kSZ: 10+/-2\% at k=1 only (2512.02954). At k=1, HACC is about 1.2-1.4 sigma below 2309.02920 and consistent with HSC-XXL. At k=3, HACC is low by a factor of about 2.1-2.3, which is about 3.1-3.4 sigma below the 2309.02920 lower error bar, even at the edge of the HACC design range. This gap is flagged as unexplained. Possible contributors: (a) HACC's minimum lies above k=3.7 h/Mpc, while FLAMINGO puts it at k\textasciitilde{}8-10 h/Mpc at z=0, so the curve shapes may differ; (b) the cosmologies and halo-mass weightings differ (2309.02920 is calibrated on haloes above \textasciitilde{}3e13 Msun, 2512.02954 on 13\textless{}log10 M500\textless{}14 at 0\textless{}z\textless{}1); (c) the HACC reference spectrum is P\_grav from the emulator, not a matched DMO run; (d) kappa\_w and M\_seed were held fixed in these HACC runs. The 2309.02920 paper itself states that it excludes some hydro simulations as having too-weak feedback.

\textbf{Finding:} Correction to the f-004 wording. Amon \& Efstathiou (2206.11794) show that suppressing the non-linear power with A\_mod\textasciitilde{}0.69 is sufficient to reconcile KiDS-1000 with Planck LCDM. It is one possible route, not the only one. On mitigation, Siegel et al. (2512.02954) find in their LSST-Y1 mock that an uninformed baryonification model removes the S8 bias but roughly doubles the S8 uncertainty (precision halved). A kSZ+eROSITA-calibrated prior recovers about half of that lost precision.

\textbf{Finding:} Parameter bounds and meanings for HACC Pk-ratio from describe\_parameters: kappa\_w is the 'AGN wind coupling / wind efficiency parameter' [2.0, 4.0]; e\_w is the 'AGN energy efficiency / wind energy fraction' [0.2, 1.0]; M\_seed is the black hole seed mass in units of 1e6 Msun [0.6, 2.0]; v\_kin is the kinetic feedback velocity in units of 1e4 km/s [0.1, 1.2]; epsilon\_kin is the kinetic feedback efficiency in units of 10 [0.02, 1.2]; omega\_m is the physical matter density, omega\_m = Omega\_m h\textasciicircum{}2 [0.12, 0.155]; sigma\_8 spans [0.7, 0.9]. The design has 110 simulations in 400 Mpc/h boxes.

\textbf{Finding:} One-at-a-time sensitivity of HACC Pk-ratio at z=0. Except for the parameter varied, all runs use the fiducial point (e\_w=0.6, v\_kin=0.65, eps\_kin=0.61, M\_seed=1.3, kappa\_w=3, omega\_m=0.14, sigma\_8=0.8). Suppression, 1-P\_hydro/P\_grav (emulator 1-sigma in brackets):

\begin{itemize}
  \item kappa\_w=2: 2.92\% (0.15\%) at k=0.997; 10.93\% (0.54\%) at k=3.007; 13.66\% (0.64\%) at k=3.714 h/Mpc.
  \item kappa\_w=4: 1.31\% (0.13\%) at k=0.997; 4.95\% (0.46\%) at k=3.007; 6.60\% (0.54\%) at k=3.714 h/Mpc.
  \item M\_seed=2.0: 2.42\% (0.14\%) at k=0.997; 9.03\% (0.49\%) at k=3.007; 11.40\% (0.57\%) at k=3.714 h/Mpc.
  \item Fiducial, from f-001: 2.17\% (0.10\%) at k=0.997, 8.03\% (0.35\%) at k=3.01, 10.2\% at k=3.71 h/Mpc.
Lower kappa\_w therefore strengthens suppression. Ranking by shift from fiducial:

  \item At k=1: lowering e\_w, v\_kin and eps\_kin together to their lower bounds (-2.4 pp, d-01 weak point) \textgreater{} kappa\_w over 2 to 4 (1.6 pp spread; +0.75 pp at kappa\_w=2, -0.86 pp at kappa\_w=4) \textgreater{} raising e\_w, v\_kin and eps\_kin together to their upper bounds (+0.33 pp) \textgreater{} M\_seed 1.3 to 2.0 (+0.25 pp).
  \item At k=3: the joint lowering (about -8.8 pp) \textgreater{} kappa\_w (5.98 pp spread; +2.9 pp at kappa\_w=2) \textgreater{} M\_seed (+1.0 pp) \textgreater{} the joint raising (+0.76 pp).
At k=3 every one of these shifts exceeds the emulator 1-sigma of 0.35-0.54\%. At k=1 the M\_seed shift and the joint-raising shift are only about 2 sigma.

\end{itemize}
\textbf{Finding:} Restatement of f-002, restricted to the sampled points only. Six in-range points have been sampled: fiducial, strong, weak, kappa\_w=2, kappa\_w=4 and M\_seed=2. At k=1 h/Mpc and z=0, their suppression ranges from -0.2\% to 2.9\%; the maximum is 2.92+/-0.15\% at kappa\_w=2. None reaches 5\% or 10\%. The most suppressed point:

\begin{itemize}
  \item is consistent with HSC-XXL (5+/-4\%) within about 0.5 sigma;
  \item falls short of eROSITA+kSZ (10+/-2\%, 2512.02954) by about 3.5 sigma.
At k=3 h/Mpc the largest sampled suppression is 10.9+/-0.5\% (kappa\_w=2). This is below the 18.4 (+2.6/-3.1)\% from 2309.02920 by about 2.4 sigma, using the lower observational error. This is a check of the sampled models only. Combined extremes, for example kappa\_w=2 together with high e\_w, v\_kin, eps\_kin and M\_seed, were not sampled, and no claim is made about the full design range.

\end{itemize}
\textbf{Completed --- Close the literature sourcing gaps and cross-check the HACC suppression amplitude against published calibrated values at k=1 and k=3 h/Mpc..} Two sourcing gaps are closed, with verbatim quotes stored in /Users/nramachandra/hep-genesis-testing/projects/ramp-combo-baryonic-pk/runs/261017\_143153/references.bib; the depth of the FLAMINGO z=0 minimum is still open. (1) 2309.02920 (key arxiv230902920), finding f-009: suppression of 0.042 (+0.012/-0.014) at k=1 h/Mpc and 0.184 (+0.026/-0.031) at k=3 h/Mpc in the fiducial model, or 0.049 and 0.179 in the alternative model; at 97.5\% credibility the impact is below 1\% for k\textless{}0.37 h/Mpc. (2) 2410.17109 (key arxiv241017109), finding f-010: the FLAMINGO fiducial minimum moves from k=6 h/Mpc at z=2 to k=10 h/Mpc at z=0. For thermal-AGN fgas variants from -10 to +4 sigma it stays at k=8-10 h/Mpc, and fgas mostly changes the amplitude. The response saturates below fgas-8sigma, and the jet model has a similar k-range. The text gives no number for the minimum depth; it is shown only in the paper's own Figs. 2, 5 and 6, so it remains open as q-08. (3) Comparison table, finding f-011, which replaces f-005. Suppression at k=1 / k=3 h/Mpc: HACC fiducial 2.2\% / 8.0\%; HACC strong 2.5\% / 8.8\%; 2309.02920 4.2\% / 18.4\%; HSC-XXL 5+/-4\% at k=1; eROSITA+kSZ 10+/-2\% at k=1. At k=1, HACC is consistent with HSC-XXL and about 1.3 sigma below 2309.02920. At k=3, HACC is low by a factor of about 2.1-2.3, or about 3.1-3.4 sigma, even at the edge of its design range. This gap is flagged as unexplained. Possible contributors: the HACC curve shape (minimum above k=3.7 h/Mpc, against 8-10 h/Mpc in FLAMINGO), the P\_grav reference, differences in cosmology and halo-mass weighting, and kappa\_w and M\_seed held fixed. No more strongly suppressed d-03 or d-05 point was available to add. (4) Correction to f-004, finding f-012: A\_mod\textasciitilde{}0.69 is sufficient to reconcile KiDS with Planck, not the only route. An uninformed baryonification model roughly doubles the S8 uncertainty, and a kSZ+eROSITA prior recovers about half of that loss. Open question q-08: the numerical depth of the FLAMINGO minimum, and whether varying kappa\_w and M\_seed lets HACC reach about 18\% suppression at k=3.

\textbf{Completed --- Find out which HACC emulator parameters control the z=0 Pk-ratio suppression, and resolve the curve beyond k=3.7 h/Mpc. This is a budget-limited replacement for d-03..} This direction is partly done. The parameter bounds, the sensitivity ranking and the 5\%/10\% check were completed. The suppression minimum is not resolved above k=3.73 h/Mpc, and the comparison plot was not made. Results are in f-013, f-014 and f-015, and the dead end is x-01.

Parameter bounds, quoted from describe\_parameters (f-013):

\begin{itemize}
  \item kappa\_w is the 'AGN wind coupling / wind efficiency parameter', range [2, 4].
  \item e\_w is the 'AGN energy efficiency / wind energy fraction', range [0.2, 1].
  \item M\_seed is in units of 1e6 Msun, range [0.6, 2].
  \item v\_kin is in units of 1e4 km/s, range [0.1, 1.2].
  \item eps\_kin is in units of 10, range [0.02, 1.2].
  \item omega\_m is Omega\_m h\textasciicircum{}2, range [0.12, 0.155].
  \item sigma\_8 has range [0.7, 0.9].
\end{itemize}
New z=0 Pk-ratio runs, each with all other parameters at the fiducial point:

\begin{itemize}
  \item kappa\_w=2: suppression 2.92+/-0.15\% at k=1, 10.93+/-0.54\% at k=3 and 13.66\% at k=3.71 h/Mpc.
  \item kappa\_w=4: 1.31+/-0.13\% at k=1 and 4.95+/-0.46\% at k=3 h/Mpc.
  \item M\_seed=2.0: 2.42+/-0.14\% at k=1 and 9.03+/-0.49\% at k=3 h/Mpc.
  \item Fiducial, for reference: 2.17\% at k=1 and 8.03\% at k=3 h/Mpc.
Lower kappa\_w strengthens suppression.

\end{itemize}
Sensitivity ranking (f-014):

\begin{itemize}
  \item At k=1: lowering e\_w, v\_kin and eps\_kin together (-2.4 pp) \textgreater{} kappa\_w (1.6 pp spread across its range) \textgreater{} raising e\_w, v\_kin and eps\_kin together (+0.33 pp) \textgreater{} M\_seed (+0.25 pp).
  \item At k=3: kappa\_w (6.0 pp spread across its range) \textgreater{} M\_seed (+1.0 pp) \textgreater{} raising e\_w, v\_kin and eps\_kin together (+0.76 pp).
  \item At k=3 every shift exceeds the emulator 1-sigma of 0.35-0.54\%. At k=1 the M\_seed and joint-raising shifts are only about 2 sigma.
\end{itemize}
Observational comparison, restricted to the six sampled points (f-015, which restates f-002):

\begin{itemize}
  \item No sampled point reaches 5\% or 10\% suppression at k=1 h/Mpc; the maximum is 2.9\% at kappa\_w=2.
  \item That point is about 0.5 sigma from HSC-XXL (5+/-4\%) and falls short of eROSITA+kSZ (10+/-2\%) by about 3.5 sigma.
  \item At k=3 h/Mpc the largest sampled value, 10.9\%, falls short of the 18.4\% from 2309.02920 by about 2.4 sigma.
  \item Combined extremes, such as kappa\_w=2 together with strong e\_w, v\_kin, eps\_kin and M\_seed, were not sampled.
\end{itemize}
Unresolved above k=3.73 h/Mpc (x-01): every curve is still deepening at k=3.73 h/Mpc, so no values at k=5 or 10 h/Mpc, no minimum and no high-k upturn were obtained. The cause is that read\_text\_file truncates the 799-bin CSVs. The comparison plot was not made because the budget ran out.

CSVs are in /Users/nramachandra/hep-genesis-testing/projects/ramp-combo-baryonic-pk/runs/261017\_143153/\{kappaw2\_z0,kappaw4\_z0,mseed2\_z0\}/Pk-ratio\_z0.0000\_prediction.csv.

Still open: q-01 (whether the k range reaches 10 h/Mpc) and q-06 (reading the CSVs beyond k=3.7 h/Mpc). q-02 and q-04 are answered for the sampled points.

\textbf{Critic, wave 2} --- continue: 7 findings stand up to the evidence, 3 questioned ('Second mitigation method (Maraio, Hall \& Taylor, 2410.12500): a theoretical-error covarian': The only evidence (tc-261017\_1-w2-d-07-73) is the title and author header of 2410.12500. The claimed specifics are not in the excerpt: ell\_max = 5000, a Euclid-like survey, bias in Omega\_m and S8 reduced to acceptable levels, and a modest increase in variance. No cite call with verbatim quotes was made (d-07 says cite was never run). The title supports only the method itself, a theoretical-error covariance.; 'Approximate comparison only. The HACC emulator gives about 0.8-0.9\% suppression at k=0.5 h': The statement 'at most 2.5\% at k=1 h/Mpc' is contradicted by f-014 in the same wave, where kappa\_w=2 gives 2.92+/-0.15\%. The conclusion that 'within its design range the emulator cannot approach' the required suppression overreaches: f-015 states that combined extremes were not sampled. The comparison is also not like-for-like. A\_mod scales P\_NL - P\_L, so the implied 1 - P/P\_DMO at k = 0.2-1 h/Mpc must be computed before it can be set against HACC ratios. None of the evidence contains that conversion.; 'Range-of-model comparison of the suppression 1-P\_hydro/P\_DMO at z=0. HACC fiducial: 2.2\% a': The k=3 conclusion, 'low by a factor of about 2.1-2.3, or about 3.1-3.4 sigma, even at the edge of the HACC design range', is superseded by f-014 and f-015. The kappa\_w=2 point gives 10.9\% at k=3, about 2.4 sigma below 18.4\% (lower error 3.1). The fiducial and strong points held kappa\_w and M\_seed fixed, so they are not at the edge of the design range. The arithmetic for the fiducial and strong points is correct. The 2309.02920 values rest on a title-only excerpt here and are carried over from f-009.).

Issues: The quoted phrases 'not as strong as inferred using KiDS data' and 'k \textasciitilde{} 0.2 h/Mpc' are attributed to the abstract, but the excerpt from tc-...-72 is cut off after the author affiliations. Only the conclusions chunk (tc-...-78) visibly supports 'stronger than favoured in recent hydrodynamic simulations'. The DES Y3 A\_mod value is still unsourced, and 2305.09827 is not cited.; The visible excerpts support the fgas range from -10 to +4 sigma and the statement that only Illustris and cosmo-OWLS 8.7 are out of reach. The excerpts do not show the minimum positions (k=6 at z=2, k=10 at z=0), the saturation below fgas-8sigma, the Bigwood et al. compatibility, or the claim that Amon \& Efstathiou is out of reach. These rest on the cite call. The depth of the FLAMINGO minimum is still open (q-08), and it is the benchmark most needed for comparing curve shapes.; The excerpts for every CSV are cut off at k of about 0.12 h/Mpc, so the k=1 and k=3 values cannot be checked against the quoted evidence. They are internally consistent with each other and with f-001. e\_w, v\_kin and eps\_kin were only moved jointly, so their individual sensitivities are unknown. omega\_m and sigma\_8 were never varied, so the claimed cosmology independence of about 1\% (FLAMINGO) has not been checked for HACC..

Stopping after wave 2: budget reached (2/2 waves used). Writing the report from the board.

\textit{Research report:}\par
\subsubsection*{Baryonic suppression of the matter power spectrum: HACC emulator predictions at z=0 and a literature assessment of AGN feedback and weak-lensing impact}
\subsubsection*{Abstract}
We used the HACC hydrodynamics emulator to predict the z=0 ratio P\_hydro/P\_grav at six parameter points, holding cosmology fixed at omega\_m=0.14 and sigma\_8=0.8.

\begin{itemize}
  \item \textbf{Suppression at k=1 h/Mpc.} It runs from an enhancement of 0.2\% at the weakest AGN point to 2.9\% suppression at the lowest wind coupling (kappa\_w=2).
  \item \textbf{Suppression at k=3 h/Mpc.} The largest sampled value is 10.9\%.
  \item \textbf{Comparison with observations.} Every sampled point is consistent with the HSC-XXL gas-fraction constraint of 5+/-4\% at k=1. All of them fall well short of the eROSITA+kSZ preference of 10+/-2\%. At k=3 they also fall short of the cluster-calibrated 18.4\% from arXiv:2309.02920.
  \item \textbf{Unresolved curve shape.} The emulator curves are still deepening at k$\approx$3.7 h/Mpc, the furthest point we could read. The suppression minimum and the high-k upturn are therefore not resolved.
\end{itemize}
From the literature, AGN feedback in groups and clusters ejects gas and drives the suppression. Lower cluster gas fractions and jet-mode feedback strengthen it. For weak lensing, reconciling KiDS or DES shear with Planck through baryons alone requires suppression stronger and reaching larger scales than current simulations predict. FLAMINGO concludes that baryons are unlikely to drive the S8 tension.

We computed no lensing observable.

\subsubsection*{Background}
The effect of galaxy formation on the total matter power spectrum was established with the OWLS simulations (arXiv:1104.1174). A library of 92 simulation power spectra provides benchmark ratios (arXiv:1906.00968). The commonly quoted shape is a 10--30\% minimum at k$\approx$5--20 h/Mpc followed by a high-k upturn. That shape was flagged in the survey as general knowledge and was not sourced from these papers during the run.

Within the texts that were read, the mechanism is as follows:

\begin{itemize}
  \item \textbf{Gas redistribution.} Energy injected by AGN redistributes gas relative to the dark matter. This suppresses small-scale power by tens of percent relative to dark-matter-only runs (arXiv:2512.02954).
  \item \textbf{Outward and inward baryon motion.} Below about 10 Mpc, AGN push baryons outward. On smaller scales, star formation concentrates baryons in halo centres (arXiv:2309.02920).
  \item \textbf{What strengthens the response in FLAMINGO} (arXiv:2410.17109):
  \item lower gas fractions in groups and clusters;
  \item collimated jets in place of thermal AGN feedback;
  \item lower stellar mass at fixed gas fraction.
  \item \textbf{Redshift dependence in FLAMINGO.} The response grows as redshift falls.
  \item The fiducial minimum moves from k=6 h/Mpc at z=2 to k=10 h/Mpc at z=0.
  \item The departure from unity begins at k$\approx$1 h/Mpc at z=2 and at k$\approx$0.1 h/Mpc at z=0.
  \item \textbf{Gas-fraction variants in FLAMINGO.} Thermal-AGN variants spanning fgas -10$\sigma$ to +4$\sigma$ keep the z=0 minimum at k=8--10 h/Mpc. These variants mainly change the amplitude of the response.
  \item The response saturates below fgas-8$\sigma$, which is an emulator extrapolation.
  \item The jet model has a similar k-range and minimum position.
  \item \textbf{Cosmology dependence and accuracy in FLAMINGO.} The cosmology dependence is about 1\% at k\textless{}10 h/Mpc, and the emulator is accurate to better than 1\% up to k=10 h/Mpc and z=3.
  \item \textbf{Depth of the FLAMINGO z=0 minimum.} The text states no number. It is shown only graphically, in Figs. 2, 5 and 6 of that paper.
\end{itemize}
Observational benchmarks at k=1 h/Mpc, for haloes with 13\textless{}log10 M500\textless{}14 at 0\textless{}z\textless{}1 (arXiv:2512.02954):

\begin{longtable}{@{}p{0.320\textwidth}p{0.320\textwidth}p{0.320\textwidth}@{}}
\toprule
\textbf{Constraint} & \textbf{Suppression at k=1 h/Mpc} & \textbf{Agreement with simulations} \\ \midrule \endhead
eROSITA+kSZ & 10+/-2\% & Only the strong-feedback runs BAHAMAS 8.0 and FLAMINGO fgas-8$\sigma$ agree \\
HSC-XXL & 5+/-4\% & --- \\
\bottomrule\end{longtable}
The same paper notes that the simulations differ from each other by more than 10\% at k=1 h/Mpc.

Grandis et al. (arXiv:2309.02920) derive cluster-based baryon-correction results:

\begin{longtable}{@{}p{0.320\textwidth}p{0.320\textwidth}p{0.320\textwidth}@{}}
\toprule
\textbf{Model} & \textbf{Suppression at k=1 h/Mpc} & \textbf{Suppression at k=3 h/Mpc} \\ \midrule \endhead
Fiducial & 0.042 (+0.012/-0.014) & 0.184 (+0.026/-0.031) \\
Alternative & 0.049 & 0.179 \\
\bottomrule\end{longtable}
At 97.5\% credibility they find an impact below 1\% for k\textless{}0.37 h/Mpc. They state that some hydrodynamical simulations have feedback that is either too strong or too weak.

For weak lensing, the relevant results are:

\begin{itemize}
  \item \textbf{KiDS-1000.} Amon \& Efstathiou (arXiv:2206.11794) show that KiDS-1000 shear, with S8 = 0.759 (+0.024/-0.021), can be reconciled with Planck $\Lambda$CDM if the non-linear power is suppressed with A\_mod$\approx$0.69. This is a sufficient route, not shown to be the only one. The suppression must extend to k$\approx$0.2 h/Mpc, outside the BAHAMAS range.
  \item \textbf{DES Y3.} Preston, Amon \& Efstathiou (arXiv:2305.09827) find that DES Y3 requires a weaker suppression than KiDS. It is still more extreme than recent hydrodynamic simulations predict. They note that underestimating baryonic effects is not the only possible interpretation.
  \item \textbf{FLAMINGO.} FLAMINGO (arXiv:2309.07959) finds that ''it is unlikely that baryon physics is driving this modification.''
  \item \textbf{Mitigation by gas calibration.} Siegel et al. (arXiv:2512.02954) use an LSST-Y1-like mock.
  \item Ignoring baryons biases S8 significantly.
  \item An uninformed baryonification model removes the bias but roughly doubles the S8 uncertainty.
  \item A kSZ+eROSITA-calibrated prior recovers about half of the lost precision.
  \item The bias is shown only graphically.
  \item \textbf{Mitigation by theoretical-error covariance.} Maraio, Hall \& Taylor (arXiv:2410.12500) replace scale cuts with a theoretical-error covariance in P(k). For a Euclid-like survey with ell\_max=5000, this reduces the bias in $\Omega$\_m and S8 to acceptable levels with a modest increase in variance. This statement rests on that paper's abstract and was not independently verified by the critic.
\end{itemize}
Other relevant tools named in the survey are BCemu (arXiv:2108.08863), the baryonic correction model (arXiv:1810.08629) and the analytic model of arXiv:2504.15633. None of these was available to run.

\subsubsection*{Results}
\paragraph{HACC emulator: fiducial, strong and weak AGN points}
The HACC Pk-ratio emulator has a 7-parameter design built from 110 simulations in 400 Mpc/h boxes. Its parameters and ranges are:

\begin{longtable}{@{}p{0.320\textwidth}p{0.320\textwidth}p{0.320\textwidth}@{}}
\toprule
\textbf{Parameter} & \textbf{Meaning} & \textbf{Range} \\ \midrule \endhead
kappa\_w & AGN wind coupling & [2, 4] \\
e\_w & AGN energy efficiency & [0.2, 1] \\
M\_seed & Black-hole seed mass, in 10\textasciicircum{}6 Msun & [0.6, 2] \\
v\_kin & Kinetic feedback velocity, in 10\textasciicircum{}4 km/s & [0.1, 1.2] \\
eps\_kin & Kinetic efficiency, in units of 10 & [0.02, 1.2] \\
omega\_m & $\Omega$\_m h$^2$ & [0.12, 0.155] \\
sigma\_8 & --- & [0.7, 0.9] \\
\bottomrule\end{longtable}
The reference spectrum is the emulator's P\_grav, not a matched dark-matter-only run.

All runs are at z=0 with omega\_m=0.14 and sigma\_8=0.8. The ''fiducial'' point is the centre of the parameter ranges, not a calibrated model. The three points are:

\begin{itemize}
  \item \textbf{Fiducial:} e\_w=0.6, v\_kin=0.65, eps\_kin=0.61.
  \item \textbf{Strong:} e\_w=1.0, v\_kin=1.2, eps\_kin=1.2.
  \item \textbf{Weak:} e\_w=0.2, v\_kin=0.1, eps\_kin=0.02.
\end{itemize}
All three use kappa\_w=3 and M\_seed=1.3.

\begin{longtable}{@{}p{0.240\textwidth}p{0.240\textwidth}p{0.240\textwidth}p{0.240\textwidth}@{}}
\toprule
\textbf{k [h/Mpc]} & \textbf{Fiducial} & \textbf{Strong} & \textbf{Weak} \\ \midrule \endhead
0.10 & 0.99968 & 0.99962 & 1.00009 \\
0.51 & 0.99197 & 0.99062 & 1.00096 \\
1.00 & 0.97832+/-0.00096 & 0.97506+/-0.00175 & 1.0023+/-0.0025 \\
2.00 & 0.95104+/-0.00229 & 0.94463+/-0.00405 & 1.0089 \\
3.0 & 0.91968+/-0.00346 & 0.91207+/-0.00596 & 1.0074 \\
3.71 & 0.8979 & 0.8904 & 1.0012+/-0.0105 \\
\bottomrule\end{longtable}
Expressed as suppression at k=1 h/Mpc, the values are 2.2\% (fiducial), 2.5\% (strong) and -0.2\% (weak). The negative value means the weak point is enhanced.

The weak point shows no suppression anywhere in the readable range. Its ratio peaks at about 1.0097+/-0.0072 near k$\approx$2.4 h/Mpc. We do not have a sourced physical explanation for this enhancement in the emulator.

Moving e\_w, v\_kin and eps\_kin together to their upper bounds changes the k=1 suppression by only about 0.3 percentage points relative to the fiducial. That shift is comparable to the emulator 1$\sigma$ at k=1, which is 0.1--0.25\%. The weak-to-fiducial difference of 2.4 pp is well resolved.

The suppression in the fiducial and strong runs is still deepening at k=3.71 h/Mpc. Their minimum therefore lies at higher k and was not measured.

Figure 8 overlays the three z=0 ratio curves (fiducial, strong, weak) produced by the emulator plotting tool:

[[FIGURE 8]]

Figure 9 shows the fiducial prediction with its emulator uncertainty band:

[[FIGURE 9]]

Warning: the board's quoted CSV excerpts stop near k$\approx$0.12 h/Mpc. Every value at k$\geq$0.5 h/Mpc is therefore taken from the run summaries and could not be checked against the quoted tool output.

The original statement that the full AGN design range spans only -0.2\% to +2.5\% at k=1 was contested. It is contradicted by the kappa\_w=2 run below. That statement is superseded by the six-point summary given later in this section.

\paragraph{One-at-a-time sensitivity: kappa\_w and M\_seed}
Each run below varies one parameter and keeps all others at the fiducial values. Suppression is 1-P\_hydro/P\_grav, with the emulator 1$\sigma$ in brackets.

\begin{longtable}{@{}p{0.240\textwidth}p{0.240\textwidth}p{0.240\textwidth}p{0.240\textwidth}@{}}
\toprule
\textbf{Run} & \textbf{k=0.997} & \textbf{k=3.007} & \textbf{k=3.714 h/Mpc} \\ \midrule \endhead
Fiducial & 2.17\% (0.10\%) & 8.03\% (0.35\%) & 10.2\% \\
kappa\_w=2 & 2.92\% (0.15\%) & 10.93\% (0.54\%) & 13.66\% (0.64\%) \\
kappa\_w=4 & 1.31\% (0.13\%) & 4.95\% (0.46\%) & 6.60\% (0.54\%) \\
M\_seed=2.0 & 2.42\% (0.14\%) & 9.03\% (0.49\%) & 11.40\% (0.57\%) \\
\bottomrule\end{longtable}
Lower wind coupling kappa\_w strengthens the suppression.

Ranking the shifts from the fiducial gives the following.

\begin{itemize}
  \item \textbf{At k=1:} the joint lowering of e\_w, v\_kin and eps\_kin (-2.4 pp) is the largest. It is followed by kappa\_w across its range (1.6 pp spread: +0.75 pp at kappa\_w=2, -0.86 pp at kappa\_w=4). Next is the joint raising of e\_w, v\_kin and eps\_kin (+0.33 pp), and last is M\_seed from 1.3 to 2.0 (+0.25 pp).
  \item \textbf{At k=3:} the joint lowering of e\_w, v\_kin and eps\_kin (about -8.8 pp) is again the largest shift sampled. It is followed by kappa\_w (5.98 pp spread; +2.9 pp at kappa\_w=2), then M\_seed (+1.0 pp), then the joint raising (+0.76 pp).
\end{itemize}
Of the parameters varied singly, kappa\_w has the largest one-at-a-time effect. However, it is not the largest shift sampled, because the joint move of the three AGN energy and kinetic parameters toward their lower bounds is larger.

At k=3 every shift exceeds the emulator 1$\sigma$ of 0.35--0.54\%. At k=1 the M\_seed shift and the joint-raising shift are only about 2$\sigma$.

e\_w, v\_kin and eps\_kin were only ever moved together, so their individual sensitivities are unknown. omega\_m and sigma\_8 were never varied. As a result, the roughly 1\% cosmology independence reported by FLAMINGO has not been checked for HACC.

For the kappa\_w and M\_seed runs, no values at k=0.1 or 0.5 h/Mpc are on the board. For the fiducial and strong points only, the large-scale suppression is about 0.03--0.04\% at k=0.1 h/Mpc and about 0.8--0.9\% at k=0.5 h/Mpc. The weak point is slightly enhanced at these k.

No comparison plot was produced for these runs.

\paragraph{Comparison with observational and calibrated benchmarks}
This is a range-of-model check, not a calibration. The cosmology and halo-mass weighting differ between HACC and the observational analyses, and HACC parameters do not map onto FLAMINGO fgas, FLAMINGO jet or BCemu parameters.

\textbf{At k=1 h/Mpc.} Across the six sampled points (fiducial, strong, weak, kappa\_w=2, kappa\_w=4, M\_seed=2), the suppression spans -0.2\% to 2.9\%. None reaches 5\% or 10\%. The most suppressed point, kappa\_w=2 at 2.92+/-0.15\%:

\begin{itemize}
  \item is about 0.5$\sigma$ from HSC-XXL;
  \item falls short of eROSITA+kSZ by about 3.5$\sigma$, using the observational error only;
  \item relative to Grandis et al., the fiducial and strong points sit about 1.2--1.4$\sigma$ below their 4.2\%.
\end{itemize}
\textbf{At k=3 h/Mpc.} The fiducial and strong HACC points give 8.0\% and 8.8\%. These are a factor of about 2.1--2.3 below the 18.4\% of arXiv:2309.02920, or about 3.1--3.4$\sigma$ against that paper's lower error bar. The largest sampled value, 10.9+/-0.5\% at kappa\_w=2, is still about 2.4$\sigma$ low.

This k=3 gap is unexplained, and the critic contested the comparison. Possible contributors are:

\begin{itemize}
  \item HACC's minimum lies beyond k=3.7 h/Mpc, while FLAMINGO's lies at 8--10 h/Mpc, so the curve shapes may differ.
  \item The cosmologies and halo-mass ranges differ. arXiv:2309.02920 is calibrated on haloes above about 3e13 Msun.
  \item HACC uses P\_grav as its reference spectrum.
  \item Combined extremes were not sampled, for example kappa\_w=2 together with high e\_w, v\_kin, eps\_kin and M\_seed.
\end{itemize}
\paragraph{Weak-lensing implications}
No C\_ell, xi\_+ or S8 shift was computed, because no tool provides them. All lensing statements in this report come from the literature.

The A\_mod-type suppression required by Amon \& Efstathiou and Preston et al. scales the whole non-linear part of P(k), with A\_mod$\approx$0.69 for KiDS. It reaches down to k$\approx$0.2 h/Mpc. The HACC fiducial and strong points suppress P(k) by less than 1\% at k$\leq$0.5 h/Mpc and by at most 2.5\% at k=1 h/Mpc.

Within the sampled range, the emulator is far from the suppression needed to reconcile shear with Planck. This comparison is approximate and was contested by the critic as a rough mapping between different parameterisations.

FLAMINGO likewise reports that the Amon--Efstathiou response lies outside the reach of the FLAMINGO emulator (arXiv:2410.17109). Both of its fgas-8$\sigma$ models, thermal and jet, fall inside the Bigwood et al. 2024 ranges, with or without kSZ.

The only sourced multipole range is that the baryon-sensitive shear extends to ell$\approx$5000 (arXiv:2410.12500). The DES Y3 numerical A\_mod was not read.

\subsubsection*{Discussion}
The emulator results agree qualitatively with the literature on two points:

\begin{itemize}
  \item Suppression grows with k through the readable range.
  \item Suppression is weak on linear scales, below 1\% at k\textless{}0.5 h/Mpc for the suppressed points. This is consistent with the \textless{}1\% bound at k\textless{}0.37 h/Mpc from arXiv:2309.02920.
\end{itemize}
Quantitatively, the sampled HACC models are mild-feedback models. At k=1 they sit at the low end of the HSC-XXL range. They cannot reproduce the strong feedback preferred by eROSITA+kSZ (arXiv:2512.02954) or the k=3 cluster-calibrated amplitude (arXiv:2309.02920).

The weak-AGN corner even shows a net enhancement of about 1\% near k$\approx$2.4 h/Mpc. The sensitivity scan suggests that the amplitude is controlled more by kappa\_w and by moving the energy and kinetic parameters toward their lower bounds than by raising them. Raising e\_w, v\_kin and eps\_kin together to their maximum adds only 0.3--0.8 pp. This saturation is reminiscent of FLAMINGO's reported saturation below fgas-8$\sigma$. However, the two parameterisations are not comparable, and this resemblance is not a demonstrated link.

There is an unresolved tension at k=3 that cannot be settled without the curve beyond k=3.7 h/Mpc. If HACC's minimum sits at k$\approx$8--10 h/Mpc like FLAMINGO's, the k=3 deficit reflects a genuinely weaker amplitude. If the shape differs, the comparison at a single k is misleading.

For weak lensing, the evidence assembled supports the FLAMINGO conclusion. Feedback consistent with X-ray gas fractions is not strong enough, nor large-scale enough, to remove the S8 tension. The KiDS and DES A\_mod solutions require suppression extending to k$\approx$0.2 h/Mpc, which is well beyond any of our sampled emulator points. The eROSITA+kSZ results favour stronger feedback than the fiducial simulations, which reopens part of this question. Even so, that paper's mitigation analysis treats baryons as a nuisance to calibrate, not as a resolution of the tension.

Dead ends:

\begin{itemize}
  \item \textbf{Reading the CSVs beyond k$\approx$3.73 h/Mpc.} The plain-text reader has no offset parameter and truncates each 799-bin CSV near k$\approx$3.73 h/Mpc. The table loader rejects the header row, and no file-writing tool was available to make a header-free copy. As a result, values at k=5 and 10 h/Mpc, the minimum and any upturn were never extracted.
  \item \textbf{Running other baryonic emulators.} No BCemu, BACCO or FLAMINGO emulator was available. Every computed number is from HACC.
\end{itemize}
\subsubsection*{Limitations}
\begin{itemize}
  \item \textbf{Truncated k range.} The k range effectively stops at 3.7 h/Mpc in our readout. The depth and position of the suppression minimum, which is a central quantity in the task, are not delivered. Whether the emulator itself reaches 10 h/Mpc is also not established.
  \item \textbf{Unverified mid- and high-k values.} The quoted evidence for every CSV ends near k$\approx$0.12 h/Mpc. Values at k$\geq$0.5 h/Mpc rest on run summaries.
  \item \textbf{Limited sampling.} Only six points were sampled, all at z=0 and at a single cosmology. There is no redshift scan, which matters because the lensing-weighted redshift is about 0.3--0.5. There is no cosmology variation, no fGas link and no combined-extreme points.
  \item \textbf{Missing literature numbers.} The FLAMINGO z=0 minimum depth and the DES Y3 A\_mod are unsourced. arXiv:2305.09827 and arXiv:2410.12500 were not added to the bibliography. The 1104.1174 claims about the minimum and upturn were never verified.
  \item \textbf{No lensing calculation.} No projected lensing calculation was possible.
  \item \textbf{Budget.} The two-wave budget was exhausted. The planned robustness and cosmology cross-checks were replaced by the limited sensitivity scan reported above.
\end{itemize}
\subsubsection*{Open directions}
\begin{enumerate}
  \item Re-read the 799-bin HACC Pk-ratio CSVs with a proper CSV reader to get the ratio at k=5 and 10 h/Mpc, the
\end{enumerate}
\end{document}
